\renewcommand{\Anno}{2002}

\renewcommand{\NumeroProblema}{1}

\renewcommand{\Testo}{Il Faraone Tetrankamon fa costruire una
piramide.  Egli fa trasportare dagli schiavi dei blocchi di pietra
della lunghezza di 10 metri, che vengono fatti scivolare su dei tronchi
cilindrici di 50 centimetri di diametro.  Supponendo che i blocchi di pietra
scorrano sui tronchi senza strisciare, e che, a loro volta, anche i
tronchi scorrano sul terreno senza strisciare, determinare di quanto
avanzano i blocchi di pietra dal momento in cui un tronco viene
inserito sotto la parte anteriore di un blocco a quello in cui lo
stesso tronco ``scivola fuori'' dalla parte posteriore del blocco.}

\renewcommand{\Risposta}{20 metri.}

\renewcommand{\Aiutino}{Mettersi nel sistema di riferimento del centro
di un tronco!  In questo sistema, se il blocco avanza di tot a
sinistra, che fa il terreno?  Dunque cosa fa il blocco rispetto al
terreno?}

\renewcommand{\SoluzioneUfficiale}{\TitoloSoluzione Mettiamoci in un
sistema di riferimento solidale con il centro di un tronco.  In tale
sistema di riferimento, quando il tronco rotola in un certo angolo
(diciamo, per fissare le idee, in senso orario) vedremo il blocco
spostarsi verso destra di un tratto $l$, e simmetricamente vedremo il 
terreno spostarsi verso sinistra di un tratto della stessa lunghezza 
$l$. Di conseguenza, mentre il blocco avanza verso destra di 10 metri 
(tutta la sua lunghezza), il terreno si sposta verso sinistra di 
altri 10 metri.

In definitiva, quindi, il blocco avanza di 20 metri rispetto al terreno.

Le seguenti figure rappresentano 5 fasi di avanzamento del blocco (si 
tenga presente che il tronco non \`{e} in scala, ed una tacca sul 
terreno permette di percepirne lo spostamento).

\begin{center}
	\psset{unit=2.4ex}
	\pspicture(-9,-4.5)(9,3)
	\newrgbcolor{sabbia}{.85 .63 .44}
	\psframe[linewidth=.6\pslinewidth](-9,-2.5)(9,2.5)
	
	\pscircle(0,0){1}
	\psarc[linewidth=.6\pslinewidth]{<-}(0,0){1.4}{145}{215}
	\psarc[linewidth=.6\pslinewidth]{<-}(0,0){1.4}{-35}{35}
	\psline[linewidth=.6\pslinewidth]{->}(-1,1.8)(1,1.8)
	\psline[linewidth=.6\pslinewidth]{<-}(-1,-1.8)(1,-1.8)
	\psframe[fillstyle=solid,fillcolor=gray](-8,1)(0,1.5)
	\psframe[fillstyle=solid,fillcolor=sabbia,linestyle=none](-8,-1.5)(8,-1)
	\psline(-8,-1)(8,-1)
	\psline(0,-1)(0,-1.3)
	\rput(0,-3.3){Fase 1}
	
	\endpspicture
	\hfill
	\pspicture(-9,-4.5)(9,3)
	\newrgbcolor{sabbia}{.85 .63 .44}
	\psframe[linewidth=.6\pslinewidth](-9,-2.5)(9,2.5)
	
	\pscircle(0,0){1}
	\psarc[linewidth=.6\pslinewidth]{<-}(0,0){1.4}{145}{215}
	\psarc[linewidth=.6\pslinewidth]{<-}(0,0){1.4}{-35}{35}
	\psline[linewidth=.6\pslinewidth]{->}(-1,1.8)(1,1.8)
	\psline[linewidth=.6\pslinewidth]{<-}(-1,-1.8)(1,-1.8)
	\psframe[fillstyle=solid,fillcolor=gray](-6,1)(2,1.5)
	\psframe[fillstyle=solid,fillcolor=sabbia,linestyle=none](-8,-1.5)(8,-1)
	\psline(-8,-1)(8,-1)
	\psline(-2,-1)(-2,-1.3)
	\rput(0,-3.3){Fase 2}
	\endpspicture
	
	\pspicture(-9,-4.5)(9,3)
	\newrgbcolor{sabbia}{.85 .63 .44}
	\psframe[linewidth=.6\pslinewidth](-9,-2.5)(9,2.5)
	
	\pscircle(0,0){1}
	\psarc[linewidth=.6\pslinewidth]{<-}(0,0){1.4}{145}{215}
	\psarc[linewidth=.6\pslinewidth]{<-}(0,0){1.4}{-35}{35}
	\psline[linewidth=.6\pslinewidth]{->}(-1,1.8)(1,1.8)
	\psline[linewidth=.6\pslinewidth]{<-}(-1,-1.8)(1,-1.8)
	\psframe[fillstyle=solid,fillcolor=gray](-4,1)(4,1.5)
	\psframe[fillstyle=solid,fillcolor=sabbia,linestyle=none](-8,-1.5)(8,-1)
	\psline(-8,-1)(8,-1)
	\psline(-4,-1)(-4,-1.3)
	\rput(0,-3.3){Fase 3}
	\endpspicture
	\hfill
	\pspicture(-9,-4.5)(9,3)
	\newrgbcolor{sabbia}{.85 .63 .44}
	\psframe[linewidth=.6\pslinewidth](-9,-2.5)(9,2.5)
	
	\pscircle(0,0){1}
	\psarc[linewidth=.6\pslinewidth]{<-}(0,0){1.4}{145}{215}
	\psarc[linewidth=.6\pslinewidth]{<-}(0,0){1.4}{-35}{35}
	\psline[linewidth=.6\pslinewidth]{->}(-1,1.8)(1,1.8)
	\psline[linewidth=.6\pslinewidth]{<-}(-1,-1.8)(1,-1.8)
	\psframe[fillstyle=solid,fillcolor=gray](-2,1)(6,1.5)
	\psframe[fillstyle=solid,fillcolor=sabbia,linestyle=none](-8,-1.5)(8,-1)
	\psline(-8,-1)(8,-1)
	\psline(-6,-1)(-6,-1.3)
	\rput(0,-3.3){Fase 4}
	\endpspicture
	
	\pspicture(-9,-4.5)(9,3)
	\newrgbcolor{sabbia}{.85 .63 .44}
	\psframe[linewidth=.6\pslinewidth](-9,-2.5)(9,2.5)
	
	\pscircle(0,0){1}
	\psarc[linewidth=.6\pslinewidth]{<-}(0,0){1.4}{145}{215}
	\psarc[linewidth=.6\pslinewidth]{<-}(0,0){1.4}{-35}{35}
	\psline[linewidth=.6\pslinewidth]{->}(-1,1.8)(1,1.8)
	\psline[linewidth=.6\pslinewidth]{<-}(-1,-1.8)(1,-1.8)
	\psframe[fillstyle=solid,fillcolor=gray](0,1)(8,1.5)
	\psframe[fillstyle=solid,fillcolor=sabbia,linestyle=none](-8,-1.5)(8,-1)
	\psline(-8,-1)(8,-1)
	\psline(-8,-1)(-8,-1.3)
	\rput(0,-3.3){Fase 5}
	\endpspicture
\end{center}}

\renewcommand{\SoluzioniAlternative}{
}

\renewcommand{\Osservazioni}{\TitoloOsservazione Il diametro del
tronco \`{e} del tutto ininfluente ai fini del calcolo
dell'avanzamento.  Quello che dipende dal diametro \`{e} il numero di
giri che fa il tronco durante l'avanzamento.  }

\renewcommand{\Commenti}{
}